\documentclass{beamer}
\usetheme{metropolis}
\usepackage{listings}
\usepackage{xcolor}

\lstset{
  basicstyle=\ttfamily\scriptsize,
  breaklines=true,
  keywordstyle=\color{blue!70!black},
  stringstyle=\color{red!60!black},
  commentstyle=\color{green!50!black},
  showstringspaces=false,
}

\title{Casino Slots}
\subtitle{Project Structure \& React Concepts in Practice}
\author{}
\date{}

\begin{document}

\maketitle

% ---------------------------------------------------------------
\section{Project Layout}

\begin{frame}{What's in the Repository?}
  \begin{itemize}
    \item \texttt{index.html} --- the single HTML page that loads the app
    \item \texttt{package.json} --- dependencies and npm scripts
    \item \texttt{vite.config.js} --- build tool configuration
    \item \texttt{src/main.jsx} --- the JavaScript entry point
    \item \texttt{src/App.jsx} --- the slot machine component
    \item \texttt{src/styles.css} --- all styling
    \item \texttt{logo.png} --- the jackpot symbol asset
  \end{itemize}
\end{frame}

\begin{frame}{The Toolchain in One Slide}
  \begin{description}
    \item[React 18] the UI library; turns components into DOM
    \item[Vite] dev server + bundler; reloads on save, compiles JSX
    \item[\texttt{@vitejs/plugin-react}] teaches Vite the JSX syntax
    \item[npm] installs packages from \texttt{package.json}
  \end{description}
  \vfill
  Three scripts: \texttt{npm run dev}, \texttt{build}, \texttt{preview}.
\end{frame}

\begin{frame}{How the App Starts}
  \begin{enumerate}
    \item Browser loads \texttt{index.html}.
    \item That file has one \texttt{<div id="root">} and a \texttt{<script>} tag pointing at \texttt{main.jsx}.
    \item \texttt{main.jsx} calls \texttt{ReactDOM.createRoot(...).render(<App />)}.
    \item React mounts \texttt{App} inside the root div and takes over.
  \end{enumerate}
\end{frame}

\begin{frame}[fragile]{The Entry Point}
\begin{lstlisting}[language=JavaScript]
import React from 'react';
import ReactDOM from 'react-dom/client';
import App from './App';
import './styles.css';

ReactDOM.createRoot(document.getElementById('root'))
  .render(
    <React.StrictMode>
      <App />
    </React.StrictMode>,
  );
\end{lstlisting}
  \texttt{StrictMode} double-invokes effects in dev to surface bugs.
\end{frame}

% ---------------------------------------------------------------
\section{Component Architecture}

\begin{frame}{Two Components, Clear Roles}
  \begin{description}
    \item[\texttt{App}] the ``smart'' parent --- owns all state, game logic, event handlers.
    \item[\texttt{Reel}] a ``dumb'' child --- receives props and renders one reel.
  \end{description}
  \vfill
  A common React pattern: push state up, keep leaves presentational.
\end{frame}

\begin{frame}{Data Flow}
  \begin{center}
    \textbf{Props flow down. Events flow up.}
  \end{center}
  \vfill
  \begin{itemize}
    \item \texttt{App} passes \texttt{symbols}, \texttt{centerIndex}, \texttt{spinning}, \texttt{highlighted} to each \texttt{Reel}.
    \item User clicks a button $\rightarrow$ handler in \texttt{App} updates state $\rightarrow$ React re-renders.
    \item \texttt{Reel} never mutates anything --- it just displays.
  \end{itemize}
\end{frame}

% ---------------------------------------------------------------
\section{React Concepts in Action}

\begin{frame}[fragile]{Concept 1: JSX}
\begin{lstlisting}[language=JavaScript]
<button
  className="spin-button"
  onClick={handleSpin}
  disabled={!canSpin}
>
  {spinning ? 'Spinning...' : 'Spin'}
</button>
\end{lstlisting}
  \begin{itemize}
    \item Braces \texttt{\{...\}} embed JavaScript inside markup.
    \item \texttt{onClick} takes a function reference, not a string.
    \item \texttt{className} (not \texttt{class}) because \texttt{class} is a reserved word.
  \end{itemize}
\end{frame}

\begin{frame}[fragile]{Concept 2: \texttt{useState} --- Reactive Memory}
\begin{lstlisting}[language=JavaScript]
const [credits, setCredits] = useState(500);
const [spinning, setSpinning] = useState(false);
const [bet, setBet] = useState(20);
\end{lstlisting}
  \begin{itemize}
    \item Each call gives a value and its setter.
    \item Calling the setter re-renders the component.
    \item The app tracks 10 pieces of state this way.
  \end{itemize}
\end{frame}

\begin{frame}[fragile]{Concept 3: \texttt{useMemo} --- Cached Computation}
\begin{lstlisting}[language=JavaScript]
const reels = useMemo(buildReels, []);

const visibleGrid = useMemo(
  () => createVisibleGrid(reels, centerIndices),
  [centerIndices, reels],
);
\end{lstlisting}
  \begin{itemize}
    \item Recomputes only when dependencies change.
    \item Empty array \texttt{[]} means ``compute once, ever.''
    \item Keeps render cheap on large/derived data.
  \end{itemize}
\end{frame}

\begin{frame}[fragile]{Concept 4: \texttt{useRef} --- Non-Reactive Storage}
\begin{lstlisting}[language=JavaScript]
const timersRef = useRef([]);

timersRef.current.push({ type: 'timeout', id });
\end{lstlisting}
  \begin{itemize}
    \item Persists across renders like state, but does \textbf{not} trigger re-renders.
    \item Perfect for timers, animation IDs, DOM handles.
    \item Here: holds every pending \texttt{setTimeout} / \texttt{requestAnimationFrame} so we can cancel them.
  \end{itemize}
\end{frame}

\begin{frame}[fragile]{Concept 5: \texttt{useEffect} --- Lifecycle \& Cleanup}
\begin{lstlisting}[language=JavaScript]
useEffect(() => {
  return () => {
    timersRef.current.forEach(({ type, id }) => {
      if (type === 'frame') cancelAnimationFrame(id);
      else clearTimeout(id);
    });
  };
}, []);
\end{lstlisting}
  \begin{itemize}
    \item Runs side effects outside the render phase.
    \item The returned function runs on unmount --- cancels timers so nothing leaks.
    \item \texttt{[]} dependency = ``mount once, clean up once.''
  \end{itemize}
\end{frame}

\begin{frame}[fragile]{Concept 6: Props and Destructuring}
\begin{lstlisting}[language=JavaScript]
function Reel({ symbols, centerIndex, spinning, highlighted }) {
  ...
}

<Reel
  symbols={reel}
  centerIndex={centerIndices[reelIndex]}
  spinning={!settledReels[reelIndex]}
  highlighted={highlightedReels[reelIndex]}
/>
\end{lstlisting}
  \begin{itemize}
    \item Parent passes data as HTML-like attributes.
    \item Child destructures the props object for readability.
  \end{itemize}
\end{frame}

\begin{frame}[fragile]{Concept 7: Lists and \texttt{key}}
\begin{lstlisting}[language=JavaScript]
{reels.map((reel, reelIndex) => (
  <Reel
    key={`reel-${reelIndex}`}
    symbols={reel}
    ...
  />
))}
\end{lstlisting}
  \begin{itemize}
    \item \texttt{map} turns an array of data into an array of elements.
    \item \texttt{key} lets React track items efficiently across renders.
  \end{itemize}
\end{frame}

\begin{frame}[fragile]{Concept 8: Conditional Rendering}
\begin{lstlisting}[language=JavaScript]
<main className={`app-shell ${lastWin ? 'win-state' : ''}`}>

{symbol === JACKPOT_SYMBOL
  ? <img src={logoUrl} alt="Jackpot logo" />
  : <span>{symbol}</span>}

{spinning ? 'Spinning...' : 'Spin'}
\end{lstlisting}
  \begin{itemize}
    \item Ternaries pick between two JSX branches.
    \item Template literals compose class names from state.
  \end{itemize}
\end{frame}

\begin{frame}[fragile]{Concept 9: Functional State Updates}
\begin{lstlisting}[language=JavaScript]
setCredits((current) => current - bet);

setSettledReels((current) =>
  current.map((v, i) => (i === reelIndex ? true : v))
);
\end{lstlisting}
  \begin{itemize}
    \item When the new state depends on the old, pass a function.
    \item Safer than \texttt{setCredits(credits - bet)} --- avoids stale closures when multiple updates queue up.
  \end{itemize}
\end{frame}

\begin{frame}{Concept 10: One-Way Data Flow}
  \begin{enumerate}
    \item State lives in \texttt{App}.
    \item Render computes JSX from that state.
    \item Events run handlers that call setters.
    \item Setters mark state dirty; React re-runs the render.
    \item The DOM is updated minimally by the reconciler.
  \end{enumerate}
  \vfill
  No two-way binding, no manual DOM manipulation.
\end{frame}

% ---------------------------------------------------------------
\section{Putting It Together}

\begin{frame}{The Spin Lifecycle}
  \begin{enumerate}
    \item User clicks \textbf{Spin} $\rightarrow$ \texttt{handleSpin()}.
    \item Pick random targets; deduct bet via \texttt{setCredits}.
    \item \texttt{scheduleSpinAnimation} loops with \texttt{requestAnimationFrame}, updating \texttt{centerIndices} each frame.
    \item Staggered \texttt{setTimeout}s mark each reel \texttt{settled}.
    \item \texttt{finishSpin} evaluates the center row, credits payout, sets the message.
    \item React re-renders; win-state class lights up the cabinet.
  \end{enumerate}
\end{frame}

\begin{frame}{Separation of Concerns}
  \begin{description}
    \item[Pure functions] \texttt{buildReels}, \texttt{createVisibleGrid}, \texttt{getLineWin} --- testable game math with no React.
    \item[Components] render JSX from props/state.
    \item[CSS] owns animation (\texttt{.reel-strip.spinning}), glow, layout.
    \item[React] glues it all together through state.
  \end{description}
\end{frame}

\begin{frame}{Takeaways}
  \begin{itemize}
    \item A tiny app still exercises the core of React: components, props, state, effects, refs, memo.
    \item Keep logic in pure functions; keep React for rendering + state.
    \item Push state up, pass data down, handle events via callbacks.
    \item Always clean up timers and subscriptions in \texttt{useEffect}.
  \end{itemize}
\end{frame}

\begin{frame}[standout]
  Questions?
\end{frame}

\end{document}
